\section{Discussion and conclusion}\label{sec:conclusion}

In this study, we developed PepALD, a chemically grounded generative framework for HELM-represented macrocyclic peptides. The model combines residue-level HELM syntax, chemistry-prior-knowledge-informed monomer representation, autoregressive latent diffusion, R-group-aware ring-bond prediction, and diffusion-adapted preference optimization. This design addresses two limitations that are particularly important for macrocyclic peptide discovery: atom-level molecular strings are long and difficult to optimize in sparse peptide data regimes, whereas purely symbolic HELM language models do not directly expose the underlying chemistry of natural and non-natural monomers. By generating each residue through a context-conditioned diffusion process in the Uni-Mol latent space and then projecting to chemically admissible HELM tokens, PepALD preserves the compactness and topology awareness of HELM while allowing the generator to exploit molecular similarity between monomers.

The empirical results support this formulation. Before target-specific optimization, PepALD outperformed HELM-GPT and PepMDLM in validity, uniqueness, novelty, diversity, and SNN profiles. The monomer-coverage analysis further showed that the gain was not merely a diffuse increase in peptide-level variability: PepALD sampled more diverse monomers than HELM-GPT, indicating broader use of the available building blocks. After permeability-oriented fine-tuning, PepALD shifted toward the high-permeability regime while retaining substantially higher validity and diversity than the permeability-optimized baselines. Finally, in two receptor-specific case studies, WP-DPO-enhanced PepALD improved docking score distributions without requiring target-specific supervised fine-tuning, while maintaining favorable predicted permeability, solubility, and internal diversity. These results suggest that a permeability-enriched PepALD prior can serve as a reusable starting point for macrocyclic peptide generation targeting new proteins.

Despite PepALD's promising performance in generating bioactive cyclic peptides, a key limitation persists in its building block generation paradigm. The current approach generates latent embeddings via diffusion and matches them to the closest pre-existing building blocks in our curated library. While this ensures high synthetic feasibility, it imposes a hard boundary on accessible chemical space, fundamentally constraining peptide diversity to the library's size and composition. To overcome this limitation, we propose extending the diffusion trajectory directly from generated latent embeddings to produce entirely novel building blocks, rather than terminating at the matching step. This next-generation framework would integrate predictive modules for synthetic feasibility and chemical stability into the diffusion network, enabling the model to learn chemical validity and synthetic accessibility during denoising. This paradigm shift would liberate PepALD from pre-defined library constraints, unlocking vastly expanded cyclic peptide chemical space and facilitating the discovery of next-generation peptide therapeutics.

In conclusion, PepALD provides a unified and extensible framework for \textit{de novo} macrocyclic peptide design. By coupling HELM grammar with chemically meaningful monomer embeddings and diffusion-based preference optimization, it generates diverse valid macrocycles and can be redirected toward permeability- and affinity-relevant objectives using reinforcement learning. These properties make PepALD a promising computational platform for exploring non-natural macrocyclic peptide space and for accelerating the discovery of cell-permeable cyclic peptides against challenging intracellular targets. We anticipate that future wet-lab experiments will further validate the ability of PepALD to generate experimentally actionable macrocyclic peptide candidates.
